Community building is critical for the success of the FCC project, because it fosters political, public, and financial support, encourages scientific collaboration and interdisciplinary innovation, promotes edu\-ca\-tion and outreach, and helps ensure the project’s long-term impact and sustainability. 
By creating an engaged, well-informed, and enthusiastic community, FCC can position itself not only as a groundbreaking scientific endeavour but also as a global, inclusive project with far-reaching benefits for humanity. 
This aspect has been given high priority ever since the beginning of the Conceptual Design Study, in 2014, in several directions:

\begin{itemize}
\item {\bf Global scientific networks}: 
The FCC project will be one of the largest scientific endeavours ever attempted and its success crucially depends on the cooperation of a broad international scientific community. 
Engaging researchers, engineers, and institutions worldwide in discussions and collaborations, early on, will ensure that 
FCC leverages the best scientific expertise, knowledge, training, and innovation from across the globe,
from the preparation of the theory and experimental tools to the actual operation of the collider and the detectors, culminating in the data analysis and interpretation.
\item {\bf Transdisciplinary collaboration}: 
The FCC project encompasses a wide range of disciplines, from particle physics theory to detector design and R\&D, from computing to engineering, from accelerator physics to machine-detector interface, etc. 
A well-established community always encourages cross-disciplinary collaboration and cross-fertilisation, enabling experts from various fields to contribute their knowledge and solve complex problems together. 
It also provides the resilience necessary to overcome setbacks, adapt to new circumstances, keep pushing the project forward to its implementation.
\item {\bf Inspiring the next generation of scientists and engineers}: 
Community engagement will help inspire the next generation to pursue careers in high-energy physics, and to build and operate the collider and the experiments. 
The FCC project must be a focal point for educational programmes, workshops, and professional training that generate interest and excitement in scientific discovery. The work done today will likely benefit future generations in ways that cannot yet be predicted. 
By creating a strong, multi-generational community, the project ensures that these long-term benefits are recognised, advocated for, and maximised over time.
\item {\bf Funding:}
Last but not least, the FCC project will require funding from multiple sources, including governments, private donors, and international organisations. 
A strong, active community can advocate for the project, raising awareness about its scientific and societal benefits. 
Demonstrating broad support and interests from the community and advocacy of the long-term scientific benefits is required to establish credibility and increase the likelihood of securing the necessary resources for the project.
\end{itemize}

As a matter of fact, the FCC feasibility study midterm review committees issued a number of recommendations to this effect. 
In particular, they recommended:

\begin{itemize}
\item to work with the scientific community, institutes, laboratories, and funding agencies to ensure support and resources for four experiments, facilitating the exploitation of the full scientific potential offered by the large investment in the FCC facility;
\item to dedicate additional human and financial resources to the project, with a resource-loaded schedule of work and clear priorities;
\item to develop the coordination and structure required to enable the theoretical progress needed to match the anticipated experimental precision of the FCC data, both at CERN (fellows, scientific associates, visitors) and by engaging collaborating institutes (including, for instance, the creation of European networks); and
\item to establish a dedicated FCC team in the research sector at CERN, with specific new positions associated, and to quantify its size and makeup in terms of seniority so that the resources required can be estimated.
\end{itemize}

Such actions have been (at least partially) anticipated a decade ago during the Conceptual Design Study (2014-2019), and intensified during the Feasibility Study (2021--2025). 
The three volumes of the FCC Conceptual Design Report~\cite{fcc-phys-cdr,fcc-ee-cdr,FCC-hhCDR}, released in January 2019, have been signed by 1364 authors, distributed in about 140 institutions, which had become members of the FCC Collaboration with the signature of Memoranda of Understanding (MoU) during this first phase. 
To comply with the recommendations of the 2021 European Strategy Update, the FCC Collaboration was reinforced by inviting more institutions to join and carry out the FCC Feasibility Study. 
One of the goals was to gather a majority of particle physicists behind the project, able to build a consensus in the community at the next European Particle Physics Strategy update that FCC is, indeed, the post-LHC collider option with the broadest scientific impact, at the intensity and the energy frontiers. 

During the five years of the Feasibility Study, the expansion of the Collaboration was pursued in two different ways.

\begin{itemize}

\item[A.] At the Collaboration level, the FCC Global Collaboration (FGC) Working Group continuously engages with current and new participants ---national institutes, laboratories and universities, as well as industry--- to encourage an expanded membership through Memoranda of Understanding. 
The FGC explores opportunities for future prospective participants, supports new participants in the application process, and assists the new participants in defining areas of collaboration, in particular for the accelerator. 
It then concludes relevant agreements to facilitate the integration process. 
It also prepares the foundations for R\&D and contributions by industry, 
as well as fostering interest in geology, geodesy, logistics, materials science, and other areas that, 
while not being at the core of CERN activities, 
are critical for the success of the FCC feasibility study and its implementation on the ground.

\item[B.] From the side of the FCC Physics, Experiments, and Detectors (PED) group, one of the six pillars of the FCC Feasibility Study, contacts are established with all high-energy physics groups in Europe and around the world, to invite them to join and contribute to the FCC project, in parallel to their current activities.
To this effect, an International Forum of National Contacts (IFNC) was created, chaired by two national contacts, with the mandate of attracting the different HEP groups, country by country. 
The quasi-complete list of European states (around 30) are represented in the IFNC, which meets regularly to reinforce the collaboration. 
The IFNC also reaches out to the rest of the world.  
The United States of America are now strongly involved in the Collaboration, following the recent Snowmass process and the positive recommendation of the U.S.\ Particle Physics Project Prioritization Panel (P5) to engage in an off-shore Higgs Factory. 
Most of the other large non-European countries are also active and represented in the IFNC (including Argentina, Brazil, Canada, Chile, India, Mexico, Pakistan, South Korea, Thailand, and Turkey) and the list continues expanding.  
\vskip .05cm
The cases of Japan and China, currently considering projects for colliders on their soil, are treated separately. 
With the goal to integrate teams from these countries in the near future, scientific exchanges with the Japanese and Chinese communities already occur in the FCC/CEPC/LC/ECFA and other national workshops.
Informal discussions with early-career Japanese physicists have also been organised.
\vskip .05cm
To build an extensive community of HEP physicists, the IFNC is also developing a finer structure, identifying institutional contacts for the participating institutes inside a country (up to 20 or 30 individuals for certain countries, such as France, Germany, Italy or the UK, and even more for the US), 
the goal being to have a vast majority of HEP institutes contributing, or on the verge of contributing, to the project by the end of the next European Strategy Update (2026).
\end{itemize}

In parallel, the PED group has also launched a call for Expressions of Interests (EoI's), in October 2024, to encourage the different institutes to get together and collaborate on innovative FCC sub-detector R\&D within the DRD Collaborations, and on detector concept studies, with the FCC integration as main objective. 
A first version of these sub-detector EoI's will be submitted as input to the European Strategy Group by March 31$^\text{st}$, 2025, and will be the basis for continuing development and consolidation during the next phase of the FCC study, until the end of 2027, with more R\&D, prototype construction, detailed simulation, test beams, etc. 
Detector concepts EoI’s will follow a similar development. 
Different combinations of sub-detectors will be integrated and tested in the common software to reach optimal performance, possibly as a function of the many physics objectives. 
Further steps are also envisioned, such as documenting, as an additional input to the European Strategy, the potential contribution of the different countries to FCC in the coming years, should the project be recommended and then approved.

The work with the scientific community progresses on a daily basis in the working groups of the PED pillar and gets a broader exposure twice a year during 
the `FCC Collaboration Weeks' in Spring (2021 at CERN; 2022 in Paris; 2023 in London; 2024 in San Francisco; 2025 in Vienna) 
and the `FCC Physics Workshops' in Winter (2022 in Liverpool; 2023 in Krakow; 2024 in Annecy; 2025 at CERN), 
allowing the community to be reinforced and new scientific collaborations to emerge. 
A strong FCC participation in the US Community Study on the Future of Particle Physics (`Snowmass 2021') acted as a decisive seed for the FCC effort in the US and was followed by three US-FCC workshops in 2023, 2024, and 2025. 
Besides, many national FCC Collaboration workshops have taken place, in essentially all large countries or regions of Europe and in the US, with active participation from the PED coordination group, yielding a significant growth of the PED part of the Collaboration in the past five years.

Several funding agencies have already reacted very positively to this evolution, decisively supported by substantive diplomatic work from the CERN management. 
The most resounding three examples are listed below.

\begin{enumerate}
\item In April 2024, the White House and CERN signed a joint statement of intent~\cite{WHCERN} saying, in particular, \emph{``Should the member states determine that FCC-ee is likely to be CERN’s next world-leading research facility, the US intends to collaborate on its construction and physics exploitation, subject to appropriate domestic approvals''.}
\item In June 2024, the CERN Council approved the Medium Term Plan (MTP) for the period 2025--2029~\cite{CERNmtp2}, including the funding of a bottom-up resource request for FCC-specific new positions.
\item In September 2024, the EU Competitiveness Report~\cite{EUCompetitivenessReport} was publicly released, with the following statements: 
\emph{``One of CERN’s most promising current projects, with significant scientific potential, is the construction of the Future Circular Collider (FCC): a 90-km ring designed initially for an electron collider and later for a hadron collider. [\ldots] 
Refinancing CERN and ensuring its continued global leadership in frontier research should be regarded as a top EU priority''.}
\end{enumerate}

Following the approval of the CERN MTP in June 2024, CERN created a dedicated FCC group in the EP department, effective since September 2024, with dedicated new positions (fellows, students, scientific associates, visitors) over the next three years. 
This sends another strong signal to the HEP community worldwide, regarding the host-lab commitment to FCC. 
Past experience has shown that even a moderately-sized host-lab group provides a significant leverage on contributions from the particle-physics community at other institutes. 
It is now anticipated that the creation of this group will be a strong incentive for CERN contract holders to reassign part of their time to the FCC project, starting with the supervision of the new recruits.
As the project moves to the next phase, many more engineers and detector physicists will be needed soon. 
The current situation is being carefully monitored and it is expected that a significant pool of these engineers and physicists will transition to FCC detector development as the  construction for the HL-LHC upgrades begins to wind down.

Assuming a positive recommendation of the CERN Council at around 2028, the FCC Collaboration expects that an FCC Committee (equivalent to the LHCC for the LHC experiments) will be formed by the CERN management and that a call for FCC detector Conceptual Design Reports will be issued soon after. 
Proto-collaborations (in particular, subsets of the FCC Collaboration, but not only) could then answer this call shortly after 2030 with full detector concept proposals. 

Finally, the coordination and structuring of the theoretical work needed to match the anticipated experimental precision of the FCC data has started during the Feasibility Study. 
Several successful mini-workshops were organised (Targets and Tools, Flavour Physics Programme, BSM Physics Programme, Higgs/Top/EW Physics Programme, Parton Shower, Phenomenology) with between 100 and 350 participants. 
The CERN/TH FCC team currently consists of three physicists, with the occasional participation of up to eight staff members and fellows, reflecting a genuine academic interest in working on the FCC physics among theorists. 
Needless to say, more dedicated resources will be needed for a worldwide organisation during the next phases of the FCC study. 